\section{Interface utilisateur et présentation des résultats}

\subsection{Choix de configuration}

L'interface permet à l'utilisateur de choisir un problème, une instance, une
stratégie de recherche et un mode d'exploration. Des limites de temps ou de
nombre de noeuds peuvent aussi être configurées, notamment pour les instances
difficiles du N-Puzzle.

\begin{lstlisting}[caption={Exemple d'une session utilisateur.}]
=== AUSSI V2 ===

Probleme :
1. Labyrinthe
2. Navigation
3. N-Puzzle
4. Jeu arithmetique

> 1

Strategie :
1. BFS
2. UCS
3. Greedy Best-First
4. A*

> 4

Mode :
1. Tree-Search
2. Graph-Search

> 2
\end{lstlisting}

\subsection{Format de sortie}

Après une exécution, \AUSSI{} affiche les informations suivantes :

\begin{itemize}
    \item Le statut de l'exécution : solution trouvée, problème insoluble,
    limite de temps ou limite de noeuds atteinte.
    \item Le chemin ou la trace de raisonnement.
    \item Le coût total de la solution.
    \item La profondeur de la solution.
    \item Le nombre de noeuds générés.
    \item Le nombre de noeuds expansés.
    \item La taille maximale de la frontière \code{OPEN}.
    \item Le temps d'exécution.
\end{itemize}

\begin{lstlisting}[caption={Exemple de résultat affiché.}]
=== Resultat ===
Statut               : SOLVED
Strategie            : A*
Mode                 : Graph-Search
Cout total           : 18
Profondeur           : 14
Noeuds generes       : 86
Noeuds expanses      : 42
Taille maximale OPEN : 21
Temps                : 0.0038 s
Chemin               : R -> R -> D -> D -> ...
\end{lstlisting}

\subsection{Visualisation}

Pour les labyrinthes et les problèmes de navigation, la solution peut être
représentée visuellement. Dans le cas du N-Puzzle, les configurations
successives peuvent être affichées. Pour le jeu arithmétique, l'interface
présente l'expression finale et les opérations qui la composent.